\begin{lemma}{Dichtheit und H�ufungspunkte}
             {DichtheitUndHaeufungspunkte}
Seien $(M, d)$ ein metrischer Raum, $x_0\in M$
ein H�ufungspunkt von $M$, sowie $D \subseteq M$ dicht in $M$.\\
Dann gilt:
\begin{enumerate}[(1)]
	\item \label{DichtheitUndHaeufungspunkte1}
	Es gibt $A \subseteq M\setminus\meng{x_0}$ abz�hlbar, sodass
	$(D\setminus\meng{x_0}) \cup A$ dicht in $M$ ist.
	\item \label{DichtheitUndHaeufungspunkte2}
	Ist $D$ abz�hlbar, so gibt es $D'' \subseteq M$ abz�hlbar und dicht in
	$M$ mit $x_0 \notin D''$.
\end{enumerate}
\end{lemma}
\begin{proof}
Da $x_0$ ein H�ufungspunkt von $M$ ist, gibt es eine gegen $x_0$ konvergente
Folge $\folge xn\N\in(M\setminus\meng{x_0})^\N$.\\Dann ist
$A:=\menge{x_n}{n\in \N} \subseteq M\setminus\meng{x_0}$ und abz�hlbar.\\
Sei $D':= \left(D \setminus \meng{x_0} \right) \cup A$.\\
Dann ist $x_0\notin D'$ und falls $D$ abz�hlbar ist, ist auch $D'$ als endliche
Vereinigung abz�hlbarer Mengen abz�hlbar.\\
Da $\folge x n \N$ gegen $x_0$ konvergiert, ist
$\menge{x_n}{n \in \N} \cup \meng{x_0}$ kompakt und damit insbesondere
abgeschlossen, also 
$\overline{\menge{x_n}{n\in\N}}\subseteq\menge{x_n}{n\in\N}\cup\meng{x_0}$.\\
Offenbar ist $x_0$ als Grenzwert von $\folge x n \N$ wegen $x_n\neq x_0$ f�r
alle $n \in \N$ ein H�ufungspunkt von $\menge{x_n}{n \in \N}$ und damit 
$\overline{\menge{x_n}{n\in\N}}\supseteq\menge{x_n}{n\in\N}\cup\meng{x_0}$.\\
Also $\overline{\menge{x_n}{n\in\N}}=\menge{x_n}{n\in\N}\cup\meng{x_0}$.\\
Es gilt weiter:
\begin{align*}\overline{D'} 
&=\overline{\left(D \setminus \meng{x_0} \right) \cup \menge{x_n}{n \in \N}}
=\overline{D \setminus \meng{x_0}} \cup \overline {\menge{x_n}{n \in \N}}\\
&= \overline{D \setminus \meng{x_0}} \cup \menge{x_n}{n \in \N}\cup\meng{x_0}\\
&\supseteq D \setminus \meng{x_0} \cup \menge{x_n}{n \in \N}\cup\meng{x_0}
= D \cup\menge{x_n}{n \in \N}\text.
\end{align*}
Da $\overline {D'}$ abgeschlossen ist, enth�lt es auch den Abschluss von
$D \cup\menge{x_n}{n \in \N}$. Somit gilt wegen der Dichtheit von $D$ in M:
\[\overline{D'} \supseteq \overline{D \cup\menge{x_n}{n \in \N}} \supseteq
\overline D = M\text.\]
Die umgekehrte Inklusion ist trivial.
\end{proof}